\section{SHADE --- ewolucja różnicowa z historyczną adaptacją} \label{sec:shade}

SHADE (\textit{Success-History based Adaptive Differential Evolution}) został zaproponowany przez Tanabego i Fukunagę \cite{Tanabe2013SHADE} i rozwija ideę JADE, zastępując mechanizm adaptacji oparty na pojedynczej parze parametrów $(\mu_{CR}, \mu_F)$ \textbf{historyczną pamięcią} o rozmiarze $H$.

\subsection{Mechanizm historycznej adaptacji}
W SHADE utrzymywane są dwie historyczne tablice o rozmiarze $H$: $M_{CR} = [M_{CR,1}, \ldots, M_{CR,H}]$ oraz $M_F = [M_{F,1}, \ldots, M_{F,H}]$, które są inicjalizowane wartościami $0{,}5$. Parametry sterujące dla $i$-tego osobnika generowane są jako:
\[
CR_i = \mathcal{N}(M_{CR, r_i}, 0{,}1), \quad F_i = \mathcal{C}(M_{F, r_i}, 0{,}1)
\]
gdzie $r_i$ jest losowym indeksem z $\{1, \ldots, H\}$, $\mathcal{N}$ oznacza rozkład normalny, a $\mathcal{C}$ --- rozkład Cauchy'ego z obcięciem do $(0, 1]$ \cite{Tanabe2013SHADE}.
Po każdym pokoleniu do tablic wpisywana jest ważona średnia Lehmera wartości parametrów z osobników, które wygenerowały lepsze potomstwo:
\[
M_{CR,k} \leftarrow \mathrm{mean}_\mathrm{WL}(S_{CR}), \quad M_{F,k} \leftarrow \mathrm{mean}_\mathrm{WL}(S_F)
\]
gdzie wagi proporcjonalne są do poprawy wartości funkcji celu \cite{Tanabe2013SHADE}. Indeks $k$ jest aktualizowany cyklicznie po każdym pokoleniu, od $1$ do $H$. Przechowywanie $H$ niezależnych wartości historycznych, zamiast jednej aktualizowanej średniej, jak w JADE, sprawia, że SHADE jest odporniejszy na przedwczesną zbieżność parametrów, szczególnie w przestrzeniach wielomodalnych \cite{Tanabe2013SHADE}.

\subsection{L-SHADE --- SHADE z liniową redukcją populacji}
W 2014 roku Tanabe i Fukunaga zaproponowali rozszerzenie SHADE o mechanizm \textbf{liniowej redukcji liczebności populacji} (ang. \textit{Linear Population Size Reduction} --- LPSR), tworząc wariant L-SHADE \cite{Tanabe2014LSHADE}. Liczebność populacji jest zmniejszana liniowo od wartości początkowej $NP_{\mathrm{init}}$ do minimalnej $NP_{\mathrm{min}}$ w miarę wyczerpywania budżetu ewaluacji:
\[
NP_G = \mathrm{round}\!\left[ \frac{NP_{\mathrm{min}} - NP_{\mathrm{init}}}{MaxFES} \cdot FES + NP_{\mathrm{init}} \right]
\]

Mechanizm ten pozwala na szeroką eksplorację przy dużej populacji na początku poszukiwań, i skupić wysiłek obliczeniowy na eksploatacji pod koniec przebiegu\cite{Tanabe2014LSHADE}. L-SHADE zajął pierwsze miejsce w konkursie CEC2014, tym samym wzmacniając pozycję tej klasy algorytmów wśród metod optymalizacji numerycznej. Mimo upływu lat badania nad ewolucją różnicową są nadal aktywnym obszarem --- współczesne kierunki rozwoju tej metaheurystyki obejmują m.in. architektury wielopopulacyjne oraz mechanizmy regeneracji elit (ang. \textit{multi-population and elites regeneration}). Cao i Luan wykazali w 2024 roku, że takie podejście pozwala skuteczniej zapobiegać przedwczesnej zbieżności oraz lepiej równoważyć procesy eksploracji i eksploatacji w globalnej optymalizacji \cite{Cao2024}.

\subsection{SHADE w kontekście optymalizacji agentów}
Opisane właściwości algorytmu SHADE czynią go dobrym narzędziem do optymalizacji wag funkcji oceny agentów gier karcianych --- automatyczna adaptacja parametrów $F$ i $CR$ eliminuje konieczność ręcznego strojenia, które w środowiskach o szumnej funkcji przystosowania jest trudne. García-Sánchez i in. \cite{sanchez2020} wykazali, że algorytmy ewolucyjne z adaptacyjnymi parametrami osiągają wyniki konkurencyjne względem algorytmów przeszukiwania drzewa gry, takich jak MCTS \cite{Swiechowski2022, Browne2012}, przy mniejszym koszcie obliczeniowym na jedną decyzję. Historyczna pamięć o długości $H$ pozwala dodatkowo zachować różnorodność parametrów przez cały przebieg eksperymentu, co ma znaczenie przy szumnych ocenach wynikających z losowości efektów gry Hearthstone \cite{hoover2020}.
